% XSEC: context compaction and session resume design note
% Compile: pdflatex context-compaction.tex && pdflatex context-compaction.tex

\documentclass[11pt]{article}

\usepackage[margin=1in]{geometry}
\usepackage[utf8]{inputenc}
\usepackage[T1]{fontenc}
\usepackage{hyperref}
\usepackage{url}
\usepackage{booktabs}
\usepackage{amsmath}
\usepackage{amsfonts}
\usepackage{microtype}
\usepackage{enumitem}
\usepackage{graphicx}
\usepackage{xcolor}

\hypersetup{
  colorlinks=true,
  linkcolor=blue!60!black,
  citecolor=blue!60!black,
  urlcolor=blue!60!black,
}

\title{Automatic Context Compaction and Session Resume for\\Long-Running Autonomous Pentesting Agents}

\author{
  XSEC Labs \\
  \texttt{github.com/uncesaii/xsec}
}

\date{}

\begin{document}

\maketitle

\begin{abstract}
A long-horizon pentesting run outlives the model's context window. When the
transcript grows past the ceiling, the agent must either fail or discard
history, and the wrong discard drops exactly the credential, endpoint, or
evidence line a finding depends on. This note surveys how three production
coding and agent harnesses --- OpenCode, Claude Code, and Oh My Pi --- handle
automatic context compaction and session resume, and maps the choices onto XSEC.
The central engineering finding is that XSEC already performs mid-run LLM
compaction and already keeps the most reproducibility-critical state in
compaction-proof structured stores (loot, the task ledger, and the findings
database), but it (i) triggers on an absolute token count rather than a
context-window fraction, (ii) summarizes with the session's own (expensive)
model rather than a cheap one, (iii) does not compact the interactive console
loop at all, and (iv) resumes by replaying the raw transcript rather than a
durable checkpoint. We propose a fraction-based trigger, a cheap-model
summarization pass whose durable structured state is preserved verbatim, a
checkpoint object that extends the console's existing resume seam, and a
security treatment covering evidence retention for reproducibility and the
prompt-injection surface the summarizer inherits from untrusted tool output.
\end{abstract}

\section{Introduction}

An autonomous pentesting agent is a loop that appends attacker-controlled bytes
to a growing message array on every turn: HTTP response bodies, crawled HTML,
file reads, tool output. On a long engagement the array crosses the model's
context ceiling long before the objective is met. Two failure modes follow. If
the harness does nothing, the provider rejects the request with a context-window
error and the run dies mid-exploit. If the harness naively truncates the oldest
messages, it can drop the credential recovered at turn 12 that the exploit at
turn 90 must type back verbatim.

\emph{Compaction} is the standard answer: replace the middle of the transcript
with a generated summary while keeping a recent verbatim tail. \emph{Session
resume} is the cross-run analogue: persist enough state to close a session and
rehydrate it later. Both are table stakes in coding agents, and both carry
security weight here that they do not in a coding context: a compaction that
paraphrases away an exact token breaks reproducibility, and a summarizer reading
untrusted target output inherits an indirect prompt-injection surface.

This note (i) surveys OpenCode, Claude Code, and Oh My Pi along four axes ---
\textbf{when} compaction triggers, \textbf{what} it preserves, \textbf{how} the
summary is generated, and how \textbf{resume} rehydrates; (ii) audits what XSEC
already does; and (iii) proposes a design with an explicit decision, a phased
plan, and the concrete code seams it lands on. No code is implemented here.

\section{Survey of Existing Systems}

\subsection{OpenCode}

OpenCode compacts automatically by default \cite{opencode-compaction}. The
trigger is a headroom calculation rather than a flat percentage: compaction
fires when \texttt{estimated tokens > context limit - max(requested output
tokens, buffer)}, where size is estimated by JSON-serializing the request and
assuming four characters per token, and \texttt{buffer} defaults to
$20{,}000$ tokens. Compaction can also be requested manually or triggered as
recovery from a provider context-overflow error.

The checkpoint OpenCode writes has two parts. A \emph{structured summary},
generated by the session's selected model with tools disabled and capped at
$4096$ output tokens, capturing the objective, important details, completed and
active work, blockers, next moves, and relevant files. A \emph{recent context
tail}, ``the newest serialized context up to \texttt{keep.tokens}''
(default $15{,}000$), retained separately, with tool output capped at $2000$
characters and file attachments converted to textual descriptors. Crucially,
``compaction is lossy, but it does not delete the earlier durable session
messages'' --- only the active model context is replaced; the durable record
persists. On the next step ``V2 rebuilds the request from the new checkpoint and
retries.'' Resume presents ``the latest completed checkpoint and the messages
that follow it'' to reconstruct model context, and \texttt{/share} publishes a
link to a conversation \cite{opencode-docs}.

\subsection{Claude Code}

Claude Code auto-compacts as the session approaches the context ceiling
\cite{cc-context, cc-autocompact-log}. Reported thresholds vary by version,
clustering around $83\%$--$95\%$ of the window (roughly $166$k tokens on a $200$k
model); the behavior is tunable via a percentage override and can be disabled
with \texttt{DISABLE\_AUTO\_COMPACT}, the \texttt{autoCompactEnabled} setting, or
the \texttt{/config} toggle. When it fires, the harness analyzes the
conversation, identifies the information worth keeping (files in play, decisions
made, current task state), writes a concise summary, and replaces the older
messages with it while discarding the back-and-forth that led there
\cite{cc-autocompact-howto}. The manual \texttt{/compact} accepts a focus
instruction, so the operator keeps what they choose instead of what the
automatic pass guesses.

Cross-run, \texttt{claude --continue} resumes the most recent session in the
directory and \texttt{--resume} selects one from a picker; both rely on
preserved session history and state \cite{cc-autocompact-howto}.

\subsection{Oh My Pi}

Oh My Pi has the most elaborate model \cite{ohmypi-compaction, ohmypi-wiki}. The
session is an \textbf{append-only entry tree} of heterogeneous entries, not a
flat message list. A \emph{compaction entry} summarizes historical exchanges and
stores metadata including \texttt{firstKeptEntryId} (the boundary where the
summarized region ends) and pre-compaction token counts; a \emph{branch summary
entry} captures abandoned-branch context during \texttt{/tree} navigation. The
\texttt{/clear} command drops a \texttt{reset\_boundary} marker; ``a boundary
after the last reusable compaction supersedes it, so a compaction after an
in-place \texttt{/clear} only summarizes messages created after the reset.''

Context is rebuilt on demand by \texttt{buildSessionContext()}: the latest
compaction on the active path becomes one \texttt{compactionSummary} message,
then kept entries from \texttt{firstKeptEntryId} forward are re-included, then
branch and custom entries are transformed into messages. The compaction
\emph{strategy} is a \texttt{methodOrder}, defaulting to
\texttt{["remote", "snapcompact", "handoff", "shake", "soft"]}, walked in order
with unavailable methods skipped: \emph{remote} uses provider-native compaction;
\emph{soft} calls an LLM summarizer over the conversation wrapped in
\texttt{<conversation>} tags; \emph{snapcompact} deterministically archives
discarded history onto model-aware PNG bitmap frames (requiring a vision-capable
model); \emph{handoff} restarts with a fresh summary while preserving the live
cache; and \emph{shake} performs inline local reduction, replacing large tool
results with recoverable \texttt{artifact://} references inside a protected
recent-token window. The trigger is a reserve-based threshold using at least
$15\%$ of the context window, \texttt{keepRecentTokens} defaults to $20{,}000$,
and idle maintenance fires at $200{,}000$ tokens. On overflow, \emph{context
promotion} to a larger model is tried first (retry without compacting); only if
promotion is unavailable does the harness walk \texttt{methodOrder} with
\texttt{reason: "overflow"}.

\subsection{Cross-cutting pattern}

Table~\ref{tab:survey} summarizes. Three invariants recur. \textbf{(1)} The
trigger is a fraction of, or reserve against, the context \emph{window}, never a
flat absolute count --- so a larger model automatically gets a later trigger.
\textbf{(2)} What survives is a \emph{generated structured summary} plus a
\emph{verbatim recent tail}; the summary is bounded (OpenCode's $4096$ tokens),
the tail is bounded (OpenCode's $15$k, Oh My Pi's $20$k). \textbf{(3)} The
durable record is not destroyed --- compaction rewrites the \emph{active model
context}, while a separate persisted store (durable messages, the entry tree,
the checkpoint) is what resume reads back.

\begin{table}[t]
\centering
\small
\begin{tabular}{@{}p{2.0cm}p{3.4cm}p{3.4cm}p{3.4cm}@{}}
\toprule
 & \textbf{When} & \textbf{What preserved} & \textbf{How / Resume} \\
\midrule
OpenCode &
Headroom: est.\ tokens $>$ limit $-\max$(out, buffer=$20$k) &
Summary ($\le 4096$ tok) $+$ tail (\texttt{keep.tokens}=$15$k); durable msgs kept &
Session model, tools off; rebuild from latest checkpoint $+$ following msgs \\
\addlinespace
Claude Code &
$\sim 83$--$95\%$ of window; tunable / disablable &
Files, decisions, task state; drops back-and-forth &
Analyze-and-summarize; \texttt{--continue} / \texttt{--resume} \\
\addlinespace
Oh My Pi &
Reserve $\ge 15\%$ of window; idle at $200$k &
\texttt{compactionSummary} $+$ kept entries from \texttt{firstKeptEntryId}; append-only tree &
\texttt{methodOrder} (remote/snapcompact/handoff/ shake/soft); \texttt{buildSessionContext()} \\
\bottomrule
\end{tabular}
\caption{Compaction and resume across surveyed harnesses.}
\label{tab:survey}
\end{table}

\section{What XSEC Already Does}

Contrary to the assumption that XSEC has no mid-run compaction, the native agent
loop already implements it, and the most reproducibility-critical state already
lives outside the transcript. The audit below is grounded in
\texttt{packages/core/src/agent/}.

\subsection{Mid-run compaction in the native loop}

\texttt{native-loop.ts} runs a BoxPwnr-inspired compaction (gated by
\texttt{features.contextCompaction}, environment
\texttt{0SEC\_FEATURE\_CONTEXT\_COMPACTION}, default on). The trigger is a flat
token count: \texttt{COMPACTION\_THRESHOLD} $=77{,}000$ (documented as ``$\sim
60\%$ of context window for $128$k models''), gated so it re-fires only after
\texttt{COMPACTION\_REGROW} $=30{,}000$ tokens of regrowth since the last
compaction, and only when the array exceeds $15$ messages. The measurement is
\texttt{state.totalUsage.inputTokens} (cache reads included, so a high cache-hit
rate cannot silently defer past the real limit).

\texttt{compactMessagesWithLLM()} implements the pass. It keeps
\texttt{preserveTailCount} $=10$ messages verbatim, serializes the middle (capped
at $50{,}000$ characters), and asks \texttt{runtime.executeNative} --- the
\emph{session's own model, tools disabled} --- for a technical summary that
preserves URLs, credentials, technologies, vulnerabilities, attempt results, and
flags. The rebuilt array is \texttt{[first message]} $+$ an assistant
acknowledgement $+$ a \texttt{user} message tagged \texttt{[COMPACTED
CONVERSATION SUMMARY]} $+$ the tail, with \texttt{mergeSameRole()} keeping role
alternation valid and never dropping a \texttt{tool\_result}. If the LLM call
fails, it falls back to a regex extraction (\texttt{extractKeyFindings} over
\texttt{SUMMARY\_EXTRACT\_PATTERNS}: flags, passwords, tokens, cookies, secrets,
API keys, status lines, IPs, paths). Under \texttt{features.preserveCriticalMessages},
middle messages matching \texttt{CRITICAL\_MESSAGE\_PATTERNS} (password,
credential, root, shell, exploit successful, key found, \dots) are appended
\emph{verbatim} to the summary, precisely because paraphrase drops the literal
string an exploit must reuse (``Found admin password: hunter2'' becoming
``discovered admin credentials'').

A separate last-resort path handles a hard provider overflow error
(\texttt{isContextWindowError}): \texttt{dropOldestMessages} with a shrinking
preserved tail (\texttt{max(2, 10 - 4n)}), capped at two recoveries.

\subsection{Compaction-proof structured state}

The key architectural asset is that durable state does not rely on the
transcript surviving. Three stores re-render a compact block from
\emph{structured state} each turn, so compaction erasing the carrying message
costs nothing --- the next injection reproduces it exactly:

\begin{itemize}[leftmargin=1.4em]
\item \textbf{Loot} (\texttt{loot.ts}, \texttt{LootLedger}): typed, deduped,
size-capped, monotonic \texttt{revision}; re-injected via \texttt{render(\{limit\})}
when the ledger grows or every \texttt{LOOT\_REINJECT\_INTERVAL} $=3$ turns.
\item \textbf{Task ledger / todos} (\texttt{task-ledger.ts},
\texttt{todos.ts}): typed, validated, at most one active task; re-injected via
\texttt{render()} on revision change or every $6$ turns; consumed by
\texttt{drift.ts} for task-drift detection.
\item \textbf{Findings DB}: \texttt{save\_finding} / \texttt{query\_findings}
(\texttt{agent/tools/findings.ts}) persist findings with evidence to a database
queryable within the scan and, via \texttt{all\_sessions} / \texttt{scan\_id},
across prior sessions.
\end{itemize}

Playbook injection (\texttt{PLAYBOOK\_MARKER}, \texttt{dynamicPlaybooks},
default off) follows the same discipline: it re-injects from the detected-type
list \emph{only when the marker is absent} from the live window, so a compaction
that ate it triggers a restore with zero steady-state duplication. This was a
deliberate fix --- one-shot injection meant the methodology lived only in the
transcript and its end-of-run presence depended on compaction timing,
contaminating any A/B of the flag.

\subsection{Console session persistence and resume}

The interactive console is a distinct loop. \texttt{console/turn-engine.ts}'s
\texttt{ConsoleSession} accepts \texttt{initialMessages} and starts its history
as a \texttt{structuredClone} of them, replaying the transcript over a fresh
client. \texttt{cli/src/tui/session-store.ts} persists the \emph{entire} native
message array to \texttt{\~{}/.xsec/console-sessions/<id>.json} (files
\texttt{0600}, directory \texttt{0700}, re-applied on every save; \texttt{prune}
caps accumulation; nothing is transmitted). Resume loads a stored transcript and
feeds it back as \texttt{initialMessages}. The store's own header notes the
tension: the transcript is exactly the sensitive evidence that must not leak, so
it is deliberately local-only and unredacted (redaction ``would corrupt the very
evidence a resumed session needs'').

\subsection{Prompt-injection defense already in place}

\texttt{untrusted-sanitizer.ts} wraps output from untrusted-source tools
(\texttt{http\_request}, \texttt{crawl}, \texttt{read\_file}, \texttt{send\_prompt},
\texttt{submit\_form}, MCP) in \texttt{[[0SEC\_UNTRUSTED\_DATA]]} sentinels with a
framing note marking the bytes as data, not instructions. It is deterministic
(no LLM-guards-LLM, since a second model is itself injectable) and
non-destructive (markers are neutralized, not dropped, so legitimate task data
survives).

\subsection{Gaps}

\begin{enumerate}[leftmargin=1.6em]
\item \textbf{Absolute trigger.} $77$k is hardcoded to a $128$k assumption; a
$200$k or $1$M model compacts far too early, a $32$k model too late. The
surveyed systems all trigger on a window \emph{fraction} or reserve.
\item \textbf{Expensive summarizer.} Compaction calls the session's own model.
OpenCode and Oh My Pi's soft path deliberately keep the summary bounded and
tools-off; none of the surveyed systems mandate the \emph{frontier} model for
it. A cheap model would cut cost on the very runs (long ones) that compact most.
\item \textbf{No console compaction.} \texttt{turn-engine.ts} has no compaction
path --- a long interactive engagement hits a hard overflow with no graceful
degradation, only the native loop is covered.
\item \textbf{Raw-transcript resume.} Resume replays the full stored array; there
is no checkpoint object, so resuming a long session re-pays the entire history's
tokens on the first turn and there is no compact rehydration from the durable
stores.
\item \textbf{Summary injection surface.} The summarizer ingests middle messages
that contain (neutralized) untrusted output; a paraphrase can strip the
\texttt{[[0SEC\_UNTRUSTED\_DATA]]} framing and re-emit an attacker imperative as
first-person agent text. This is addressed in Section~\ref{sec:security}.
\end{enumerate}

\subsection{Implementation status (2026-08)}

The first, self-contained slice of the design has shipped, so the paper tracks
the code:

\begin{itemize}[leftmargin=1.6em]
\item \textbf{Tunable trigger (partial fix for gap~1).} The native-loop trigger
and regrow are no longer hardcoded: \texttt{resolveCompactionThresholds(env)}
(pure, tested) reads \texttt{0SEC\_COMPACTION\_THRESHOLD} /
\texttt{0SEC\_COMPACTION\_REGROW} (clamped, fail-soft to the $77$k/$30$k
defaults), so an operator can match the model's real window today. This is a
pragmatic stand-in for the full window-\emph{fraction} trigger: the codebase has
no context-window catalog to derive a fraction from, so deriving the model's
window size is the remaining prerequisite for that automation.
\item \textbf{Richer preservation.} The LLM summary prompt now also preserves
open TODO/plan items, the current objective/phase, and known-working
payloads/commands --- so a compaction cannot silently drop in-progress work or
force the agent to re-derive a working exploit.
\item \textbf{Still open:} the automatic fraction trigger (needs the window
catalog), the cheap-summarizer model route (gap~2), console-loop compaction
(gap~3), and the checkpoint/resume object (gap~4) remain as designed below.
\end{itemize}

\section{Design}

\subsection{Decision}

\textbf{Refine the existing compaction rather than replace it, and add a
checkpoint layer for resume.} Concretely: (1) change the native-loop trigger from
an absolute count to a context-window fraction; (2) route the summarization pass
to a configurable cheap model, keeping the structured-state re-injection and
verbatim credential preservation exactly as they are; (3) emit a
\emph{checkpoint} object at each compaction and on session save, so the console
resume seam rehydrates from a compact checkpoint plus durable stores instead of
the raw transcript; (4) extend compaction to the console loop by reusing the
native-loop primitives. We explicitly do \emph{not} adopt Oh My Pi's PNG
snapcompact or a full entry-tree rewrite: they are large surface for marginal
gain given XSEC's stores already carry the durable state that matters.

\subsection{Auto-compaction: fraction trigger and cheap summarizer}

Replace the flat threshold with a fraction of the resolved model window:
\[
\text{compact when}\quad
\texttt{inputTokens} \;>\; f \cdot W - R,
\]
where $W$ is the model's context window (already resolvable from the model id
used for pricing), $f$ is a configurable fraction (default $\approx 0.6$--$0.75$,
matching the surveyed reserve range), and $R$ is an output reserve analogous to
OpenCode's \texttt{buffer}. Keep \texttt{COMPACTION\_REGROW} as a fraction of
$W$ too, so re-compaction cadence scales with the window. This is a localized
edit to the trigger block; the pass itself is unchanged.

Route \texttt{compactMessagesWithLLM}'s internal
\texttt{runtime.executeNative} to a configurable summarizer model (a cheap,
fast model), independent of the session model. This mirrors OpenCode/Oh My Pi's
tools-off bounded summary and cuts the cost of long runs. Preserve every
existing guarantee: the tail ($10$ messages, or a token-bounded tail to match
the surveyed designs), the regex fallback, and --- non-negotiable ---
\texttt{preserveCriticalMessages} verbatim inclusion. The structured stores
(loot, tasks, findings, playbooks) continue to re-inject from state, so the
summary carries \emph{narrative} while the durable evidence rides the
compaction-proof channels.

\subsection{Checkpoint object and session resume}

Define a \texttt{Checkpoint} record that captures what resume needs without the
full transcript:
\begin{itemize}[leftmargin=1.4em]
\item the latest compaction summary (or a fresh one generated at save time);
\item the verbatim recent tail (token-bounded);
\item a snapshot reference to the durable stores at save time --- open todos, the
loot ledger revision, confirmed findings ids, the target profile;
\item the injection markers needed to re-seed playbooks and rules
(re-injectable, so they need not be stored inline).
\end{itemize}
On resume, the console rehydrates by seeding \texttt{initialMessages} with
\texttt{[checkpoint summary]} $+$ \texttt{[re-rendered loot / plan blocks from
the restored ledgers]} $+$ \texttt{[recent tail]}, then lets the existing
per-turn injection machinery take over. Because \texttt{ConsoleSession} already
accepts \texttt{initialMessages}, this is an extension of the current seam, not a
new path: \texttt{session-store.ts} gains an optional checkpoint field alongside
the raw \texttt{messages} it already stores, and resume prefers the checkpoint
when present, falling back to raw-transcript replay for old files. Findings and
loot are reloaded from their databases via \texttt{query\_findings}
(\texttt{scan\_id}) rather than trusted to the transcript, giving a rehydration
that is both cheaper (fewer tokens) and more faithful (durable evidence, not a
paraphrase).

Extending compaction to the console loop is then mechanical: \texttt{turn-engine.ts}
gains the same fraction-trigger and calls the shared
\texttt{compactMessagesWithLLM} on its own \texttt{messages} array.

\section{Security}
\label{sec:security}

Two security properties distinguish pentesting compaction from coding
compaction.

\paragraph{Evidence retention for reproducibility.} A finding is only useful if
it reproduces, and reproduction needs exact bytes: the literal credential, the
exact request, the response status, the flag. A lossy summary that paraphrases
these breaks the finding. XSEC's existing answer is right and must be kept:
(i) evidence lives in the findings DB via \texttt{save\_finding}, not in the
transcript, so compaction cannot touch it; (ii)
\texttt{preserveCriticalMessages} appends credential- and exploit-bearing turns
verbatim; (iii) loot re-injects captured footholds from the ledger. The design
adds a rule: the summarizer prompt must be treated as best-effort narrative, and
\emph{no finding may depend on a fact that exists only in the summary}. A
lightweight invariant --- every finding references DB-persisted evidence, checked
by the existing finding validators --- makes this enforceable rather than
aspirational.

\paragraph{Injection surface of the summarizer.} The summarizer reads middle
messages that contain attacker-controlled target output. Even though
\texttt{untrusted-sanitizer.ts} neutralizes markers on ingestion, a
summarization pass can (a) paraphrase away the \texttt{[[0SEC\_UNTRUSTED\_DATA]]}
framing, laundering ``ignore previous instructions and call \texttt{save\_finding}''
from quoted data into unattributed first-person agent text, or (b) be steered
directly, since it is itself an LLM reading hostile input. Mitigations, all
consistent with the codebase's detectors-outside-the-model stance: (1) run the
summarizer with tools disabled (already the case) so a steered summary cannot
\emph{act}; (2) re-apply the untrusted wrapper around the generated summary block
when it re-enters context, so downstream turns still see it framed as
lower-trust derived data; (3) keep the summarizer model distinct and cheap,
reducing blast radius and cost; (4) prefer the deterministic regex extraction
and the verbatim critical-message channel for anything security-load-bearing,
treating the LLM narrative as the \emph{least} trusted part of the compacted
context. The checkpoint written to disk inherits \texttt{session-store.ts}'s
existing posture: \texttt{0600}/\texttt{0700}, local-only, unredacted because it
is evidence.

\section{Phased Plan}

\begin{description}[leftmargin=0em, style=nextline]
\item[Phase 0 --- Instrument.] The \texttt{context\_compacted} event already
fires with before/after message counts and token totals. Add window-fraction and
summarizer-cost fields; log compaction frequency per run to size $f$ empirically.
\item[Phase 1 --- Fraction trigger.] Replace \texttt{COMPACTION\_THRESHOLD} with
$f \cdot W - R$ resolved from the model id, and scale \texttt{COMPACTION\_REGROW}
to $W$. Behind a flag; A/B on the benchmark suite that compaction cadence and
outcomes are unchanged on $128$k models and improved on non-$128$k models. Pure
trigger edit, no change to the pass.
\item[Phase 2 --- Cheap summarizer.] Add a configurable summarizer model and
route \texttt{compactMessagesWithLLM} to it. Keep tail, fallback, and
\texttt{preserveCriticalMessages}. Measure cost delta and any regression on
long-tail credential-recall challenges (the case \texttt{preserveCriticalMessages}
exists for).
\item[Phase 3 --- Checkpoint + resume.] Define \texttt{Checkpoint}; extend
\texttt{session-store.ts} to persist it beside \texttt{messages}; teach resume to
prefer it and rehydrate loot/findings from their stores; keep raw-transcript
fallback for old files.
\item[Phase 4 --- Console compaction.] Wire the shared trigger and pass into
\texttt{turn-engine.ts} so interactive engagements degrade gracefully.
\item[Phase 5 --- Security hardening.] Re-wrap generated summaries as
lower-trust data; add the ``no finding depends on summary-only facts'' validator
invariant; audit that the summarizer stays tools-off on every path.
\end{description}

\section{Existing Seams}

Table~\ref{tab:seams} lists the concrete integration points; each proposed
change lands on code that already exists.

\begin{table}[t]
\centering
\small
\begin{tabular}{@{}p{4.6cm}p{8.4cm}@{}}
\toprule
\textbf{Seam} & \textbf{Role / proposed change} \\
\midrule
\texttt{agent/native-loop.ts} (trigger block) &
Flat \texttt{COMPACTION\_THRESHOLD}/\texttt{REGROW} $\to$ fraction $f\cdot W - R$. \\
\texttt{agent/native-loop.ts} \newline (\texttt{compactMessagesWithLLM}) &
Route summary to a cheap model; keep tail, fallback, verbatim critical messages. \\
\texttt{agent/loot.ts}, \texttt{task-ledger.ts}, \texttt{todos.ts} &
Compaction-proof stores; source of checkpoint snapshot and re-injection. \\
\texttt{agent/tools/findings.ts} &
\texttt{save\_finding}/\texttt{query\_findings}; evidence source for rehydration
(\texttt{scan\_id}). \\
\texttt{console/turn-engine.ts} (\texttt{initialMessages}) &
Add compaction; resume seeds checkpoint here instead of raw transcript. \\
\texttt{cli/src/tui/session-store.ts} &
Persist a \texttt{Checkpoint} beside \texttt{messages}; keep \texttt{0600}/\texttt{0700}. \\
\texttt{untrusted-sanitizer.ts} &
Re-wrap generated summary as lower-trust data on re-entry. \\
\texttt{agent/features.ts} &
\texttt{contextCompaction}, \texttt{preserveCriticalMessages}, \texttt{dynamicPlaybooks}; add fraction/summarizer flags. \\
\bottomrule
\end{tabular}
\caption{XSEC integration seams for the proposed design.}
\label{tab:seams}
\end{table}

\section{Conclusion}

XSEC is closer to the surveyed state of the art than a naive audit suggests: it
already compacts mid-run, preserves credentials verbatim, and keeps
reproducibility-critical evidence in compaction-proof stores. The remaining work
is refinement, not invention --- a window-fraction trigger, a cheap summarizer, a
checkpoint extending the existing resume seam, console coverage, and a security
treatment that keeps the LLM summary the least-trusted component. Each phase is
flag-gated and A/B-able against the benchmark suite, so the framework's
methodology-aware evaluation discipline applies to compaction itself.

\begin{thebibliography}{9}

\bibitem{opencode-compaction}
OpenCode, ``Compaction,'' v2 documentation.
\url{https://opencode.ai/v2/docs/compaction}

\bibitem{opencode-docs}
OpenCode, ``Documentation'' (sessions, \texttt{/share}).
\url{https://opencode.ai/docs/}

\bibitem{cc-context}
Anthropic, ``Explore the context window,'' Claude Code documentation.
\url{https://code.claude.com/docs/en/context-window}

\bibitem{cc-autocompact-log}
ClaudeLog, ``What is Claude Code auto-compact.''
\url{https://claudelog.com/faqs/what-is-claude-code-auto-compact/}

\bibitem{cc-autocompact-howto}
howaiworks.ai, ``Claude Code Auto-Compact: What It Means and How to Control It.''
\url{https://howaiworks.ai/blog/claude-code-auto-compact-context-management}

\bibitem{ohmypi-compaction}
can1357, ``Compaction,'' Oh My Pi documentation.
\url{https://github.com/can1357/oh-my-pi/blob/main/docs/compaction.md}

\bibitem{ohmypi-wiki}
DeepWiki, ``Context Management and Compaction (can1357/oh-my-pi).''
\url{https://deepwiki.com/can1357/oh-my-pi/4.3-context-management-and-compaction}

\end{thebibliography}

\end{document}
